% 02_cosmology_dark_sector.tex
% Ch.2: Galactic rotation, RAR — MOND from Axiom 4

\chapter{Cosmology: Galactic Rotation from Axiom 4}
\label{ch:cosmology}

The most dramatic application of scale invariance closes the
``dark matter problem.''  The MOND acceleration scale
$a_0 \approx 1.2 \times 10^{-10}$~m/s$^2$ is first \textit{derived}:
\begin{equation}
    a_0 = \frac{c \cdot H_\infty}{2\pi}
    \approx 1.07 \times 10^{-10}\;\text{m/s}^2
    \label{eq:a0_derived}
\end{equation}
where $H_\infty$ is the asymptotic Hubble parameter from the derivation
chain (10.7\% below the empirical value).

The lattice mutual inductance $\eta$ at galactic scale is governed
by the \textit{same} saturation operator as particle confinement:
\begin{equation}
    \eta_{eff}(\dot\gamma) \;=\; \eta_0 \cdot
    \sqrt{1 - \left(\frac{\dot\gamma}{\dot\gamma_{yield}}\right)^{\!2}}
    \label{eq:galactic_inductance}
\end{equation}
When local Newtonian acceleration $g_N$ serves as the saturation
amplitude (with $a_0$ as the yield limit), the drag contribution
becomes:
\begin{equation}
    \boxed{%
    g_{eff} \;=\; g_N + \sqrt{g_N \cdot a_0}\;
    \sqrt{1 - \frac{g_N}{a_0}}
    }
    \label{eq:saturation_mond}
\end{equation}

\section{Asymptotic Limits}

\begin{itemize}
    \item $g_N \ll a_0$: $g_{eff} \to \sqrt{g_N \cdot a_0}$
          \quad (deep MOND --- flat rotation curves)
    \item $g_N \ge a_0$: $g_{eff} = g_N$
          \quad (Newtonian --- saturated lattice, no drag)
\end{itemize}

\section{Physical Mechanism}

In the outer galaxy where $g_N < a_0$,
the vacuum lattice is in its linear (unsaturated) regime.  Its mutual
inductance $\eta_{eff} \approx \eta_0$ resists differential rotation,
producing an additional radial acceleration that mimics a dark matter
halo.  In the inner galaxy where $g_N > a_0$, the same saturation
operator that drives particle confinement drives
$\eta_{eff} \to 0$~--- the lattice saturates and cannot support
transverse inductive drag.  Orbits become conservative and Keplerian.

\section{Validation: NGC~3198}

The saturation model produces a flat
rotation velocity of $\sim$159~km/s versus the observed 150~km/s
($\sim$5\% agreement), with a flatness ratio $v(25\,\text{kpc}) /
v(10\,\text{kpc}) = 0.961$~--- identical to the McGaugh empirical
interpolation.  Both the deep-MOND asymptote and the Newtonian limit
are reproduced exactly, with zero adjustable parameters.

\begin{figure}[H]
    \centering
    \includegraphics[width=0.95\textwidth]{galactic_rotation_curve.pdf}
    \caption{\textbf{Left:} NGC~3198 rotation curve. Newtonian
    prediction (dashed) falls off as $1/\sqrt{r}$; the AVE Axiom~4
    saturation (solid) produces flat curves matching observations
    (green dots). \textbf{Right:} Radial Acceleration Relation
    for three galaxies (NGC~3198, Milky Way, DDO~154). The 1:1
    line (no dark matter) and deep-MOND asymptote are shown.}
    \label{fig:galactic_rotation}
\end{figure}
